\documentclass[a4paper,12pt]{extarticle}
\input{assessments/lecture04/quiz-style.tex}
\begin{document}
\quizheader{1 · C, GCC и GDB}
\input{assessments/lecture04/common-questions.tex}
\newpage
\section*{Вариант 1 · вопросы 7–10}
\begin{enumerate}
\setcounter{enumi}{6}
\setlength{\itemsep}{6mm}
\question{Вы изменили \texttt{factorial.c}, но не запускали GCC повторно. Что выполнит \texttt{./factorial}?}
  {Автоматически соберёт изменённый исходник.}
  {Запустит прежнюю собранную программу.}
  {Откроет исходник в редакторе.}
  {Запустит GDB.}
\question{Зачем при сборке для лабораторной указан флаг \texttt{-g}?}
  {Чтобы добавить сведения, помогающие отладке в GDB.}
  {Чтобы запустить программу после сборки.}
  {Чтобы сохранить вывод программы в файл.}
  {Чтобы исправить все ошибки в коде.}
\question{Что делает GDB, когда исполнение программы остановлено на точке останова?}
  {Позволяет исследовать её состояние, в том числе значения переменных.}
  {Сразу удаляет исполняемый файл.}
  {Превращает исходный файл в PDF.}
  {Всегда перезапускает компьютер.}
\question{Какое значение должен вернуть корректный расчёт \texttt{factorial(0)}?}
  {0}{1}{-1}{Произойдёт деление на нуль.}
\end{enumerate}
\answergrid
\end{document}
